\section{Complete ODD model specification}
\label{app:odd}

This appendix supplies the self-contained Overview, Design concepts and Details specification of the model following \citet{grimm2020}. It complements the concise methodological account with implementation-level equations, initialization rules and parameter ranges.

\subsection{Overview}

\paragraph{Purpose and patterns.} The purpose is to compare how copied source attributes alter selection of an experimental false-news source and subsequent false-repost exposure under reach and recommendation policies. The model is expected to reproduce directional, not empirical time-series, patterns: greater visibility cannot reduce the set of initially informed agents; a strategy has no direct effect when the source is unreachable and undiscoverable; and longer-lived evidence permits effects to accumulate.

\paragraph{Entities and state variables.} Entity types are user agents, four source objects, endorsement events and one scenario manager. User state contains an identifier, a common vector of 14 preference weights (13 source attributes plus word of mouth, WOM), known sources, directed contacts, timestamped endorsements, pending WOM evidence and the most recent selection score. Source state contains an identifier, category, reach, 13 binary attribute distributions, objective false-publication probability, active scenario attributes, publication-state history and a cumulative set of reposter identifiers. Scenario state contains source, target, inclusive interval and copied attribute groups.

\paragraph{Scale.} The principal population is 500 users, four source objects and 400 dimensionless periods. Attribute levels are binary. The baseline directed network has five outgoing contacts per user. One publication state is generated by each source in each period.

\paragraph{Process schedule.} Initialization loads and validates the workbook, creates sources, creates 500 copies of the aggregate user prototype, samples source visibility, constructs contacts and inserts deterministic period-$-1$ endorsements. Each period then executes scenario changes, source publication draws, pending-label delivery, user activity and choice, reporting, and contact recommendation in that order. Recommendations can affect choice no earlier than the next period plus label delay.

\subsection{Design concepts}

\paragraph{Basic principles and emergence.} Bounded source choice is driven by accumulated attribute-level evidence and contact information. Population shares, exposure and unique reach emerge from repeated local choices; no aggregate trajectory is imposed.

\paragraph{Adaptation, objectives and learning.} Users probabilistically choose sources with larger aggregate scores but have no explicit utility maximization or strategic foresight. Choice creates fresh attribute endorsements, so later evaluations are path dependent. Source strategy is exogenous and scheduled; the source does not optimize from feedback.

\paragraph{Sensing and interaction.} Users sense only known sources, their endorsement histories and the strongest source recommendation supplied by directed contacts. They do not sense global popularity, network structure or cumulative false-news prevalence. A recommendation can reveal an unknown source.

\paragraph{Heterogeneity.} Preference parameters are homogeneous because the only available user record contains survey means. Visibility, network position, activity, sampled attribute experiences, publication exposure, choices and memories generate state heterogeneity. The model does not preserve respondent-level preference distributions or correlations.

\paragraph{Stochasticity.} Random draws determine initial visibility, random-network contacts or small-world rewiring, activity, source attribute levels, source choice, source publication state, label coverage and classification errors. A single seeded pseudorandom stream is used within a run. Common seed values improve treatment--control comparability but separate random streams are not used; changes that alter the number or order of draws can therefore weaken common-random-number alignment. Pairing is retained as a design match and this limitation is disclosed.

\paragraph{Collectives, observation and prediction.} Contacts form a directed graph but users do not form explicit groups. The reporter records per-period selections, publication states and cumulative unique reposters. Predictions are conditional comparisons within the encoded model, not forecasts of X.

\subsection{Details}

\paragraph{Initialization.} For each source, initial visibility is sampled independently using category reach when reach filtering is active. For every known source and source attribute, the more probable binary level is selected deterministically. If the levels are exactly tied, the first level is retained because the implementation updates the maximum only on a strict increase. With the common weight $m_a$ for attribute $a$, binary level 1 produces $-m_a/2$ and level 2 produces $m_a$. Deterministic initialization supplies a common prior-familiarity anchor without consuming treatment-dependent random draws; it removes initial attribute uncertainty and is therefore a modelling limitation. These events are dated $-1$ and treated as age zero at the first decision. The random fixed-outdegree network samples distinct non-self contacts. The alternative directed small-world network begins with forward ring neighbours and rewires each edge independently with probability 0.1; collisions are resampled, so outdegree remains fixed while indegree varies.

\paragraph{Attribute sampling and endorsement.} During interaction, a binary source level $L_{sat}\in\{1,2\}$ is sampled from the active source distribution. It creates signed value
\begin{equation}
v_{isat}=\begin{cases}-m_a/2,&L_{sat}=1,\\m_a,&L_{sat}=2.\end{cases}
\end{equation}
The score transform, memory weighting and choice normalization are given in Section~3.2.

\paragraph{Publication behaviour.} At the start of each period, source $s$ draws $Z_{st}\sim\mathrm{Bernoulli}(p_s)$. This state is independent of the current attribute realization and is shared by every user selecting the source in that period. Explicit $p_s$ therefore prevents a credibility-copying strategy from changing objective veracity. Legacy workbooks without a source-behaviour sheet instead derive $p_s$ from the current low-credibility probability; that rule is retained only for backward compatibility and is not used by the study workbook.

\paragraph{WOM and labels.} Each friend exposes its selected source and current score. For duplicate source recommendations the receiver retains the maximum score, then selects the source with the maximum retained value; exact maxima retain source-identifier iteration order. Given coverage $c$, a label is sampled independently for that receiver--recommendation event. A false state is classified false with sensitivity $s_e$; a true state is classified true with specificity $s_p$. The observed class selects a configured effect $g\in\{-1,0,1\}$. Non-zero evidence $g\lambda W$ is dated $t+1+d$, where $W$ is common because user preferences are homogeneous. Coverage failure creates no WOM endorsement, although the recommended source has already become known.

\paragraph{Scenario operation.} On its first active period, a strategy copies specified attribute distributions from traditional media to the experimental source. Several named groups can be combined. On the period after an inclusive end date, original target attributes are restored. Stored endorsements are not rewritten.

\paragraph{Output calculation.} Every active user with at least one known source contributes at most one decision. A selection counts as false when the source-period $Z_{st}=1$. Unique reach is the number of distinct users who have ever selected each source through the indicated horizon. Formal denominators are given in Section~3.4.

\begin{table*}[ht]
\centering
\caption{Principal parameters and uncertainty ranges. A dash means the value is not varied in that experiment.}
\label{tab:parameters}
\small
\begin{tabular}{p{3.4cm}p{2.1cm}p{3.2cm}p{6.2cm}}
\toprule
Parameter & Principal value & Revision range/levels & Interpretation \\
\midrule
Users; periods; source objects & 500; 400; 4 & fixed & Population, dimensionless horizon and category representatives \\
Choice base $B$ & 1.2 & 1.10--1.40 & Exponential transform of signed endorsements \\
Memory specification & exponential $h=0.33$ & hard window 25 or $h=0.33$--5 & Alternative specifications in each run: a hard cutoff, or infinite hard memory with exponential decay \\
Contacts; friend fraction & 5; 1 & contacts 3--10 & Outdegree is $\lfloor\texttt{CONTACTS}\times\texttt{FRIENDS}\rfloor$ \\
Activity probability & 1 & 0.25, 0.50, 1 & Probability that an agent makes its available decision \\
Topology; rewiring & random; 0.10 & random or directed small-world & Fixed outdegree, no self-links/duplicates; variable indegree \\
WOM receiver scale $\lambda$ & 0.5 & 0.25, 0.5, 1 & Multiplier applied to the common receiver WOM preference \\
Label delay $d$ & 0 & 0, 5, 25 & Additional periods beyond next-period delivery \\
Coverage $c$ & 1 & 0.5, 0.8, 1 & Probability that a recommendation receives a label \\
Sensitivity; specificity & 1; 1 & 0.75--1; 0.90--1 & Classification of false and true source-period states \\
Source-attribute contrast & 1 & 0.75--1.25 & Scales binary probabilities away from 0.5, clipped to $[0,1]$ \\
Experimental-source reach & 14.7\% & 0--14.7\% & Probability of initial direct visibility \\
Traditional false-publication probability & .092 & [.05,.15] & Objective source-period probability \\
Unknown-source false-publication probability & .206 & [.15,.30] & Objective source-period probability \\
Experimental-source false-publication probability & .667 & [.55,.80] & Objective source-period probability \\
Mixed-source false-publication probability & .416 & [.30,.55] & Objective source-period probability \\
\bottomrule
\end{tabular}
\end{table*}

\begin{table*}[ht]
\centering
\caption{Survey-informed binary source-attribute inputs. Cells report $P(\text{level 2})$; $P(\text{level 1})=1-P(\text{level 2})$. Attribute levels are publication-level stochastic realizations from source-level distributions.}
\label{tab:source-attributes}
\small
\begin{tabular}{lrrrr}
\toprule
Attribute & Traditional & Unknown & Experimental & Mixed \\
\midrule
Entertainment & .884 & .899 & .667 & .820 \\
Geographic proximity & .860 & .824 & .640 & .764 \\
Audiovisual content & .949 & .891 & .773 & .865 \\
Links & .908 & .807 & .453 & .708 \\
Political proximity & .814 & .828 & .413 & .652 \\
Hashtags & .775 & .777 & .800 & .843 \\
Sensationalism & .611 & .672 & .853 & .787 \\
Source credibility & .908 & .794 & .333 & .584 \\
Information quality & .874 & .853 & .427 & .618 \\
Positive emotion & .179 & .122 & .093 & .067 \\
Social proximity & .862 & .849 & .693 & .753 \\
Simple language & .237 & .429 & .693 & .584 \\
Negative emotion & .029 & .021 & .240 & .090 \\
\bottomrule
\end{tabular}
\end{table*}

